\documentclass{report}

\usepackage{amsmath,amssymb,amsthm}
\usepackage{hyperref}
\usepackage[capitalize]{cleveref}

% Theorem environments
\newtheorem{theorem}{Theorem}
\newtheorem{definition}{Definition}
\newtheorem{axiom}{Axiom}
\newtheorem{corollary}{Corollary}
\newtheorem{remark}{Remark}

% Custom macros for leanblueprint
\providecommand{\home}[1]{}
\providecommand{\github}[1]{}
\providecommand{\dochome}[1]{}
\providecommand{\lean}[1]{}
\providecommand{\leanok}{}
\providecommand{\uses}[1]{}

\title{ISAR: Invariant Scalable Fractal Addressable Representations}
\author{Fabian Schindler}
\date{\today}

\begin{document}
\maketitle

\begin{abstract}
This monograph presents the complete mathematical foundations of the ISAR (Invariant Scalable Fractal Addressable Representation) framework, a formally verified computational substrate implemented in the Lean 4 proof assistant. We detail: (1) the core discrete ISAR combinator calculus, proving confluence and unique normal forms; (2) a category-theoretic terminality proof showing that observational quotients factor uniquely through the invariant layer; (3) relative set-theoretic interpretations in Hereditarily Finite sets; (4) dialect views and partial evaluation via Futamura projections; and (5) continuous completions of metric address spaces and the Universal Approximation Theorem.
\end{abstract}

\chapter{Introduction and Foundational Philosophy}

Modern software engineering, programming language design, and formal verification are plagued by a profound fragmentation. Disjoint computational models—such as the untyped lambda calculus, term rewriting systems (TRS) and e-graphs, compiler intermediate representations (IR), arithmetic and quantity calculi, and virtual machine bytecodes—are treated as separate worlds. They occupy distinct namespaces, employ incompatible metatheories, and rely on ad-hoc, pairwise compiler pipelines that leak implementation details and fail to preserve observational boundaries.

This thesis argues that this fragmentation is unnecessary. We show that a broad class of closed formal systems can be represented through a single, quotient-mediated semantic kernel without loss of operational behavior. This kernel is not a new ``universal language,'' but rather a view-independent substrate: different programming languages and formalisms become different \textit{decoders} over the same invariant structure.

\section{Primal Axioms and Ontological Closure}

The ISAR (Invariant Scalable Fractal Addressable Representations) framework starts from the minimal assumptions of existence and distinction, leading to a mathematically closed computational substrate.

\begin{axiom}[Existence]
\leanok
The universe of discourse is non-empty: $\Omega \neq \emptyset$.
\end{axiom}

\begin{axiom}[Identity]
\leanok
The initial state is a singular, indistinguishable point (the Void, denoted $\emptyset$). Distinction begins when this Void is partitioned.
\end{axiom}

\begin{axiom}[Self-Description]
\leanok
A system exists only if it can describe its own existence. This self-description requirement forces the transition from a singleton to a pair: $P = (\emptyset, \emptyset)$.
\end{axiom}

\begin{axiom}[Asymmetry]
\leanok
The act of description creates a fundamental distinction between the \textit{Observer} and the \textit{Observed}. In set-theoretic terms, this corresponds to the Kuratowski pairing:
$$(x, y) = \{\{x\}, \{x, y\}\}$$
This asymmetry, denoted $\Delta$, serves as the operational fuel driving directed rewriting steps and computations.
\end{axiom}

\section{The 4-Carrier Minimality Proof by Exclusion}

The computational substrate is spanned by four primitive carriers: Identity ($I$), Rewrite ($R$), Adjacency ($A$), and State ($S$). We prove that exactly four independent carriers are necessary and sufficient to form an admissible computational substrate:

\begin{enumerate}
    \item \textbf{Identity ($C_I$)}: Represents the invariant operational equivalence. Without $C_I$, the substrate lacks a notion of stable normal forms or equivalence-preserving transformations, causing all state transitions to dissolve into noise.
    \item \textbf{Rewrite ($C_R$)}: Represents the directed, non-invertible selection/choice. Without $C_R$, transitions are symmetric, leading to static, reversible systems incapable of representing one-way computational steps.
    \item \textbf{State ($C_S$)}: Represents repeatable observables and state-space closure. Under self-application, the system must remain closed. Without $C_S$ to copy and distribute terms over contexts, the system collapses into a strictly linear, non-duplicating fragment that is sub-universal (incapable of simulating arbitrary Turing-complete computation).
    \item \textbf{Adjacency ($C_A$)}: Represents causal connectivity and topological interaction. It maps the relations between parameters. Without $C_A$, the carriers $C_I, C_R, C_S$ cannot be composed into ordered causal chains; the system collapses to an unordered bag of static values.
\end{enumerate}

No carrier can serve dual roles without violating single-carrier localization contexts:
\begin{itemize}
    \item If $C_I = C_S$, there is no proper quotient structure (observational equivalence collapses to identity).
    \item If $C_R = C_A$, causality becomes mutable, leading to non-deterministic, time-dependent topological rewrites that violate confluence.
    \item If $C_I = C_R$, the invariant becomes strategy-dependent, causing evaluation order to determine the semantic outcome.
\end{itemize}
Thus, the minimal size of the carrier set $C$ is exactly 4.

\section{Sizing Hierarchies: The Genesis of Nibbles and Bytes}

The token hierarchy of digital computing (bits, nibbles, bytes) is typically justified by legacy hardware constraints (e.g., ASCII encoding, memory alignment on the IBM System/360 in 1964). In the ISAR framework, this hierarchy is derived constructively from first principles of directed relation mapping:

\begin{itemize}
    \item \textbf{1 Bit}: Distinguishes existence from absence (atom vs. void).
    \item \textbf{2 Bits}: Identifies one of the four active roles ($C_I, C_R, C_A, C_S$) in the minimal carrier set.
    \item \textbf{4 Bits (One Nibble)}: Represents a directed pair of names (an input pointer to a carrier role and an output pointer from a carrier role). This is the minimum structure needed to express \textit{adjacency} (a relation between two roles).
    \item \textbf{8 Bits (One Byte)}: Represents a pair of pairs—a source context and a target context—which is the minimal faithful representation of a \textit{directed rewrite step} (a morphism between two carrier contexts).
\end{itemize}
Thus, early hardware engineers arrived at the 8-bit byte because it is the minimal structural representation of an asymmetric rewrite step ($\Delta$) with a parity check.

\section{Occam and the Description-Length Functional}

For any substrate $M$, we define the description-length marginal likelihood:
$$ p(D \mid M) = \int p(D \mid \theta, M) p(\theta \mid M) d\theta $$
where $\theta$ represents the parameter space (matrices and rules) and $D$ represents a reference computation. Given the 4-carrier limit, any substrate with $n > 4$ carriers introduces redundant gauge flexibility that increases description length. The terminal ISAR kernel minimizes this description length, acting as the unique Minimum Description Length (MDL) substrate.
\chapter{Core Combinatory Calculus}

In this chapter, we formalize the syntax and reduction semantics of the ISAR combinatory logic kernel. The core ISAR calculus is a substrate-independent combinator system that supports S-K combinator reductions, parallel reductions, and normalization.

\section{Syntax and Reduction Semantics}

The syntax consists of standard S-K combinators and the specialized ISAR terms.

\begin{definition}[S-K Term Syntax]
\lean{ISAR.SK}
The inductive type of basic S-K combinator expressions is defined by:
\begin{align*}
t, u \in \text{SK} &::= \mathbf{S} \mid \mathbf{K} \mid t\ u
\end{align*}
\end{definition}

\begin{definition}[ISAR Term Syntax]
\lean{ISAR.ITerm}
The core term syntax of the ISAR kernel extends basic combinators with norms, constants, and custom operations:
\begin{align*}
t, u \in \text{ITerm} &::= \mathbf{norm} \mid \mathbf{s_s} \mid \mathbf{konst} \mid \mathbf{app}\ t\ u
\end{align*}
\end{definition}

We define the reduction rules for both systems.

\begin{definition}[SKStep Reduction]
\lean{ISAR.SKStep}
The reduction relation on SK terms represents the standard combinatory logic reduction rules:
\begin{itemize}
    \item $\mathbf{K}\ x\ y \to x$
    \item $\mathbf{S}\ x\ y\ z \to (x\ z)(y\ z)$
\end{itemize}
\end{definition}

\begin{definition}[IStep Reduction]
\lean{ISAR.IStep}
The single-step reduction relation on ISAR terms defines the operational semantics of norms and constants.
\end{definition}

\begin{theorem}[Confluence of IRed]
\lean{ISAR.IRed_confluence}
\uses{ISAR.ITerm, ISAR.IStep}
The reduction relation $\text{IRed}$ (the reflexive transitive closure of $\text{IStep}$) satisfies the Church-Rosser confluence property:
$$\forall t, u_1, u_2, \quad t \to^* u_1 \land t \to^* u_2 \implies \exists v, \quad u_1 \to^* v \land u_2 \to^* v$$
\end{theorem}

\begin{theorem}[Unique Normal Forms]
\lean{ISAR.isar_fragment_unique_normal_forms}
\uses{ISAR.IRed_confluence}
If a term $t$ reduces to two normal forms $n_1$ and $n_2$, then $n_1 = n_2$.
\end{theorem}

\section{The Invariant Layer}

Using the uniqueness of normal forms, we define the quotient space representing behavioral equivalence classes of terms.

\begin{definition}[ISKSubtype]
\lean{ISAR.ISKSubtype}
\uses{ISAR.ITerm}
The subtype of ISAR terms belonging to the S-K fragment:
$$\text{ISKSubtype} := \{ t : \text{ITerm} \mid \text{ISKTerm } t \}$$
\end{definition}

\begin{definition}[Operational Equivalence]
\lean{ISAR.OperEq}
\uses{ISAR.ISKSubtype}
Two terms are operationally equivalent if they behave identically under reduction:
$$t \sim_{op} u \iff \text{OperEq } t\ u$$
\end{definition}

\begin{definition}[Invariant Layer]
\lean{ISAR.InvariantLayer}
\uses{ISAR.OperEq}
The quotient type of \texttt{ISKSubtype} modulo \texttt{OperEq}:
$$\text{InvariantLayer} := \text{Quotient } (\text{operEqSetoid})$$
\end{definition}

\begin{definition}[Invariant Layer Application]
\lean{ISAR.InvariantLayer.app}
\uses{ISAR.InvariantLayer}
The application operator lifted to the quotient space:
$$\text{app} : \text{InvariantLayer} \to \text{InvariantLayer} \to \text{InvariantLayer}$$
\end{definition}

\begin{definition}[Canonical Representative]
\lean{ISAR.InvariantLayer.canonical_rep}
\uses{ISAR.InvariantLayer}
A noncomputable function mapping each equivalence class to its canonical representative in \texttt{ISKSubtype}.
\end{definition}

\chapter{Category-Theoretic Interface and Terminality}

The ISAR kernel provides a category-theoretic interpretation of representation-free substrates and semantic views. We formalize this using the category of semantic kernels and morphisms, and prove that the canonical invariant quotient space is a terminal object in this category. This universal property of terminality establishes the mathematical foundation of view-independence.

\section{The Category of Semantic Kernels}

The definitions and structures in this section are mechanically verified in \url{src/ISAR/KernelCategory.lean} and \url{src/ISAR/DialectKernel.lean}. We define the category $\mathcal{K}$ of admissible semantic kernels over the substrate.

\begin{definition}[Semantic Kernel]
\lean{ISAR.Kernel}
An admissible semantic kernel $K \in \mathcal{K}$ consists of:
\begin{enumerate}
    \item A carrier type $\text{Carrier}$.
    \item A view mapping $\text{view\_of} : \text{ISKSubtype} \to \text{Carrier}$.
    \item An observational equivalence relation $\approx$ on $\text{Carrier}$.
    \item A soundness property: operational equivalence in the substrate implies view equivalence:
    $$\forall t, u, \quad t \sim_{op} u \implies \text{view\_of } t \approx \text{view\_of } u$$
    \item A decoding/reconstruction mapping $\text{decode} : \text{Carrier} \to \text{ISKSubtype}$.
    \item Coherence axioms:
    \begin{itemize}
        \item $\text{decode}(\text{view\_of } t) \sim_{op} t$
        \item $\text{view\_of}(\text{decode } c) \approx c$
        \item $c_1 \approx c_2 \implies \text{decode } c_1 \sim_{op} \text{decode } c_2$
    \end{itemize}
\end{enumerate}
\end{definition}

\begin{definition}[Canonical ISAR Kernel]
\lean{ISAR.ISAR_Kernel}
\uses{ISAR.Kernel}
The identity kernel mapping the substrate directly to itself is the canonical presentation.
\end{definition}

\begin{definition}[Computable ISAR Kernel]
\lean{ISAR.ComputableISAR_Kernel}
\uses{ISAR.Kernel}
The computable kernel parametrized by normalization fuel.
\end{definition}

\begin{definition}[Optimal Computable Kernel]
\lean{ISAR.ComputableISAR_Kernel_Optimal}
\uses{ISAR.ComputableISAR_Kernel}
The computable kernel instantiated with optimal fuel computed from the term size for linearly-typed terms.
\end{definition}

\begin{definition}[Kernel Morphism]
\lean{ISAR.KernelHom}
\uses{ISAR.Kernel}
A morphism $f : K_1 \to K_2$ in the category $\mathcal{K}$ is a carrier mapping $\text{hom} : K_1.\text{Carrier} \to K_2.\text{Carrier}$ that:
\begin{itemize}
    \item Preserves the view representations: $\forall t \in \text{ISKSubtype}, \quad K_2.\text{view\_of } t \approx_2 \text{hom}(K_1.\text{view\_of } t)$
    \item Respects carrier equivalence: $\forall c_1, c_2, \quad c_1 \approx_1 c_2 \implies \text{hom}(c_1) \approx_2 \text{hom}(c_2)$
\end{itemize}
\end{definition}

\begin{definition}[Canonical Homomorphism]
\lean{ISAR.canonical_hom}
\uses{ISAR.KernelHom, ISAR.ISAR_Kernel}
For any kernel $K$, the canonical homomorphism into $\text{ISAR\_Kernel}$ is defined by the decoding mapping.
\end{definition}

\section{Universal Factorization and Terminality}

The uniqueness and terminality theorems are formalized and proved in \url{src/ISAR/KernelCategory.lean}. We prove that the canonical representation-free kernel is terminal, meaning that any admissible view can be uniquely decoded into it.

\begin{theorem}[Morphism Uniqueness]
\lean{ISAR.morphism_uniqueness}
\uses{ISAR.KernelHom, ISAR.canonical_hom}
Any structure-preserving morphism $f : K \to \text{ISAR\_Kernel}$ is observationally equivalent to the canonical decoding morphism:
$$\forall c \in K.\text{Carrier}, \quad f(c) \sim_{op} \text{decode } c$$
\end{theorem}

\begin{theorem}[ISAR Kernel Terminality]
\lean{ISAR.ISAR_Kernel_terminal}
\uses{ISAR.morphism_uniqueness, ISAR.ISAR_Kernel}
The quotient of the morphism space $\text{KernelHom } K\ \text{ISAR\_Kernel}$ modulo observational equivalence is a singleton:
$$\exists! [f] \in \text{KernelHom } K\ \text{ISAR\_Kernel} / \sim_{op}$$
Thus, $\text{ISAR\_Kernel}$ is a terminal object in the category $\mathcal{K}$ of semantic kernels.
\end{theorem}

This category-theoretic result is not a naming convention: it is an algebraic theorem stating that the Invariant Layer is the unique (up to isomorphism) cofinal quotient that retains only the information visible to observational equivalence. Any other view is mathematically guaranteed to factor uniquely through it, separating the dialect-specific representation from the coordinate-free invariant semantics.

\section{The Dialect Subsystem}

The Dialect structures in this section are verified in \url{src/ISAR/DialectKernel.lean}. A Dialect is an alternative abstract view over the substrate.

\begin{definition}[Dialect]
\lean{ISAR.Dialect}
A Dialect specifies:
\begin{enumerate}
    \item A type of dialect objects.
    \item A type of observations and observational equivalence.
    \item An evaluation mapping, an encoder to substrate terms, and a decoder from the substrate.
    \item A coherence law: evaluation commutes with encoding and decoding.
\end{enumerate}
\end{definition}

\chapter{Symbolic Views and Dialects}

The ISAR architecture supports multiple symbolic views or ``dialects'' over the representation-free substrate. Each view compiles to the substrate and decodes back, satisfying a preservation law. In this chapter, we formalize several dialects: the Lambda calculus fragment, ZFC hereditarily finite sets, term rewriting systems, bytecode compilation, and Barker's Iota combinator.

\section{The Lambda Calculus Fragment}

The bracket abstraction compiler and simulation properties in this section are mechanically verified in \url{src/ISAR/LambdaFragment.lean}. We define the standard $\lambda$-calculus with de Bruijn indices.

\begin{definition}[Lambda Terms]
\lean{ISAR.LTerm}
Lambda terms with variables (represented by natural number de Bruijn indices), applications, and abstractions:
\begin{align*}
t, u \in \text{LTerm} &::= \mathbf{var}\ n \mid \mathbf{app}\ t\ u \mid \mathbf{lam}\ t
\end{align*}
\end{definition}

\begin{definition}[Lambda Step]
\lean{ISAR.LStep}
Single-step $\beta$-reduction on \texttt{LTerm} using de Bruijn shifts and substitutions.
\end{definition}

\begin{definition}[Lambda Compiler]
\lean{ISAR.compile}
\uses{ISAR.LTerm, ISAR.ITerm}
Compiles lambda terms into the ISAR substrate via bracket abstraction:
$$\text{compile} : \text{LTerm} \to \text{ITerm}$$
\end{definition}

\begin{theorem}[Compiler Simulation]
\lean{ISAR.compile_simulates_step, ISAR.compile_simulates_red}
\uses{ISAR.compile, ISAR.LStep}
The compiler faithfully simulates lambda reductions in the ISAR kernel:
$$\forall t, u \in \text{LTerm}, \quad t \to_\beta u \implies \text{compile } t \to^* \text{compile } u$$
\end{theorem}

\section{Hereditarily Finite Sets and ZFC}

The set-theoretic encoding and relative ZFC interpretation verified in this section are mechanically verified in \url{src/ISAR/HFSet.lean}, \url{src/ISAR/HFSetEncoding.lean}, and \url{src/ISAR/HFSetSemantics.lean}. This relative interpretation is strictly limited to the Hereditarily Finite (HF) set fragment (ZFC without the axiom of infinity). We formalize HF Sets, their encoding, and show they satisfy set-theoretic axioms.

\begin{definition}[HF Sets]
\lean{HF}
The type of hereditarily finite sets, built inductively from the empty set:
$$x, y \in \text{HF} ::= \emptyset \mid x \cup \{y\}$$
\end{definition}

\begin{definition}[Extensional Equality]
\lean{ExtEq}
\uses{HF}
Set-theoretic extensional equivalence ($x =_{ext} y$) defined via Ackermann's bijective encoding `toNat`:
$$x =_{ext} y \iff \text{ExtEq } x\ y$$
\end{definition}

\begin{definition}[HF Set Encoding]
\lean{ISAR.HF_encode}
\uses{HF, ISAR.InvariantLayer}
Encodes HF sets into substrate invariant layers:
$$\text{HF\_encode} : \text{HF} \to \text{InvariantLayer}$$
\end{definition}

\begin{theorem}[HF Encoding Correctness]
\lean{ISAR.decode_layer_HF_encode, ISAR.HF_encode_decode_layer}
\uses{ISAR.HF_encode}
The encoding is sound and invertible:
$$\forall c \in \text{HF}, \quad \text{decode\_layer}(\text{HF\_encode } c) = c$$
\end{theorem}

We package HF Sets as a semantic kernel.

\begin{definition}[HF Set Kernel]
\lean{ISAR.HF_Kernel}
\uses{ISAR.Kernel, HF}
The set-theoretic admissible semantic kernel HF\_Kernel.
\end{definition}

\begin{theorem}[ZFC Interpretation Factorization]
\lean{ISAR.HF_Kernel_factorization}
\uses{ISAR.HF_Kernel, ISAR.ISAR_Kernel_terminal}
The HF set kernel factors uniquely through \texttt{ISAR\_Kernel}:
$$\forall f : \text{KernelHom HF\_Kernel ISAR\_Kernel}, \quad f(c) \sim_{op} \text{encode\_raw } c$$
This faithful factorization proves that the ZFC HF set fragment faithfully interprets into the ISAR substrate, satisfying empty set, pairing, and union axioms.

Syntactically distinct terms representing the same set (e.g., $\{a, b\}$ and $\{b, a\}$) project to the exact same element in the `InvariantLayer` quotient, demonstrating extensional equivalence.
\end{theorem}

\section{Term Rewriting, Bytecode, and Iota Views}

The SKI TRS view, stack bytecode compiler, and Barker's iota dialect are verified in \url{src/ISAR/TRSView.lean}, \url{src/ISAR/BytecodeView.lean}, and \url{src/ISAR/IotaView.lean} respectively. We define dialects for other standard computational models.

\begin{definition}[SKI TRS Dialect]
\lean{ISAR.TRS_Dialect}
\uses{ISAR.Dialect}
The term rewriting dialect modeling pure SKI combinator terms \texttt{TTerm}.
\end{definition}

\begin{definition}[Bytecode Dialect]
\lean{ISAR.Bytecode_Dialect}
\uses{ISAR.Dialect}
The dialect representing stack-based virtual machine instruction bytecode \texttt{Instruction} compiled and decompiled.
\end{definition}

\begin{definition}[Barker's Iota Dialect]
\lean{ISAR.Iota_Dialect}
\uses{ISAR.Dialect}
The dialect modeling the single-combinator \texttt{IotaTerm} utilizing Barker's universal combinator $\iota = \lambda x. x\ S\ K$.
\end{definition}

In Barker's iota combinator view, iterating iota-application produces the closed orbit:
$$ I \to A \to K \to S \to X \to I $$
where:
\begin{itemize}
    \item $I$ = identity ($I x = x$)
    \item $A$ = flip constant ($A x y = y$)
    \item $K$ = constant ($K x y = x$)
    \item $S$ = substitution ($S f g x = f x (g x)$)
    \item $X$ = self-application kernel
\end{itemize}
These four distinct active combinators map directly onto the four matrices ($I, R, A, S$) of the ISAR substrate, proving the ontological completeness of the 4-carrier representation.

\chapter{Continuous Approximation and Representation Space}

In this chapter, we bridge the discrete combinatory logic of the ISAR substrate with continuous analysis and matrix spaces. We formalize dimensional physical quantities, matrix representations, expressive completeness, the tensor semantics, and the continuous universal approximation theorem showing that every continuous function can be uniquely represented as a limit address.

\section{Physical Quantities and Matrix Representations}

We define dimensional quantities and 4x4 matrix algebra.

\begin{definition}[Quantities and Uncertainty]
\lean{ISAR.Quantity}
A quantity represents a physical parameter with:
\begin{enumerate}
    \item A value (rational scalar).
    \item A physical dimension (exponent exponents).
    \item An epistemic uncertainty expression \texttt{Uncertainty} and \texttt{Correlation}.
\end{enumerate}
Covariance propagation propagates exact metric-epistemic covariance equations under addition and multiplication.
\end{definition}

\begin{definition}[Matrix4 Carriers]
\lean{ISAR.Matrix4}
The concrete integer 4x4 matrix representation used for basis combinator signatures. The primitive matrices are:
\begin{gather*}
I_1 = \begin{pmatrix} 1 & 0 & 0 & 0 \\ 0 & 0 & 0 & 0 \\ 0 & 0 & 1 & 0 \\ 0 & 0 & 0 & 0 \end{pmatrix}, \quad
R_1 = \begin{pmatrix} 1 & 0 & 0 & 0 \\ 0 & 0 & 0 & 0 \\ 0 & 1 & 0 & 0 \\ 0 & 0 & 1 & 0 \end{pmatrix}, \\
A_1 = \begin{pmatrix} 0 & 0 & 0 & 0 \\ 1 & 0 & 0 & 0 \\ 0 & 1 & 0 & 0 \\ 0 & 0 & 0 & 0 \end{pmatrix}, \quad
S_1 = \begin{pmatrix} 1 & 1 & 0 & 0 \\ 0 & 1 & 0 & 0 \\ 0 & 0 & 1 & 0 \\ 0 & 0 & 0 & 1 \end{pmatrix}
\end{gather*}
\end{definition}

\begin{theorem}[Idempotency and Nilpotency]
\lean{ISAR.I1_idempotent, ISAR.K1_nilpotent}
\uses{ISAR.Matrix4}
The identity matrix $I_1$ is idempotent:
$$I_1^2 = I_1$$
The core rewrite operator $K_1 = I_1 \cdot R_1 \cdot A_1 \cdot S_1$ is nilpotent:
$$K_1^2 = \mathbf{0}$$
\end{theorem}

\begin{theorem}[Gauge Similarity]
\lean{ISAR.K1_K2_gauge_equiv}
\uses{ISAR.Matrix4}
The two alternative matrix representations $K_1$ and $K_2$ are conjugate (similar) via the lower-triangular gauge transformation matrix $P$:
$$P \cdot K_1 \cdot P^{-1} = K_2$$
proving that representation details are gauge shadows, while the terminal quotient state is invariant.
\end{theorem}

\begin{definition}[Matrix View Map]
\lean{ISAR.term_signature_val}
\uses{ISAR.Matrix4, ISAR.ITerm}
The structural homomorphism mapping each ISAR term to its concrete 4x4 integer matrix signature:
$$\text{term\_signature\_val} : \text{ITerm} \to \text{Matrix4}$$
\end{definition}

\begin{theorem}[Expressive Completeness]
\lean{ISAR.isk_expressive_completeness, ISAR.basis_expressive_completeness}
\uses{ISAR.term_signature_val}
Every matrix generated by the basis combinator algebra is the image of some term:
$$\forall M \in \text{ISKAlgebra}, \quad \exists t \in \text{ISKSubtype}, \quad \text{term\_signature\_val } t = M$$
\end{theorem}

\section{The Categorical Matrix Bridge}

We package the matrix representations as an admissible semantic kernel.

\begin{definition}[Matrix Kernel]
\lean{ISAR.MatrixKernel}
\uses{ISAR.Kernel, ISAR.term_signature_val}
The semantic kernel MatrixKernel whose carrier is \texttt{Matrix4} and whose view map is \texttt{term\_signature\_val}.
\end{definition}

\begin{theorem}[Nilpotent Collapse]
\lean{ISAR.nilpotent_collapse}
\uses{ISAR.MatrixKernel}
Applying \texttt{konst} to itself maps to the zero matrix, showing the nilpotent collapse of the basis algebra:
$$\text{kernelMatrixView}(\mathbf{K}\ \mathbf{K}) = \mathbf{0}$$
\end{theorem}

\begin{theorem}[Matrix Terminality]
\lean{ISAR.matrix_kernel_terminality}
\uses{ISAR.MatrixKernel, ISAR.ISAR_Kernel_terminal}
Every structure-preserving morphism from \texttt{MatrixKernel} to \texttt{ISAR\_Kernel} is observationally equivalent to the canonical decoding morphism.
\end{theorem}

\begin{theorem}[Unreachability of R and A]
\lean{ISAR.kernel_view_R_unreachable}
\uses{ISAR.MatrixKernel}
The structural homomorphism \texttt{kernelMatrixView} never produces the rotation matrix $R_1$ or adjacency matrix $A_1$ from a pure ISK term.
\end{theorem}

\section{Tensor Denotation Semantics}

We bridge the discrete combinatory logic of the ISAR substrate with continuous analysis and matrix spaces by formalizing the denotational mapping of symbolic ISAR terms into an abstract rank-4 tensor space, where reductions map to extensional equality.

\begin{definition}[Tensor Space Primitives]
\lean{ISAR.ExtTensor}
The abstract tensor space $\mathcal{T}$ is characterized by five primitive operators:
\begin{itemize}
    \item Identity: \texttt{t\_norm}
    \item Constant: \texttt{t\_konst}
    \item Duplicator: \texttt{t\_dup}
    \item Swapper: \texttt{t\_swap}
    \item Composition: \texttt{t\_comp}
\end{itemize}
\end{definition}

\begin{definition}[Derived S in Tensor Space]
\lean{ISAR.t_sₛ}
\uses{ISAR.ExtTensor}
The distributive combinator $\mathbf{S}$ in the tensor space is constructively derived using swapper, composition, and duplicator operators:
\begin{align*}
\text{t\_s}_{\text{s}} &\triangleq \text{t\_app}\Big(\text{t\_app}\big(\text{t\_comp}, \text{t\_app}(\text{t\_comp}, \text{t\_dup})\big), \\
&\quad \text{t\_app}\big(\text{t\_app}(\text{t\_swap}, P), \text{t\_norm}\big)\Big)
\end{align*}
where:
\[ P \triangleq \text{t\_app}\big(\text{t\_app}(\text{t\_comp}, \text{t\_comp}), \text{t\_app}(\text{t\_app}(\text{t\_comp}, \text{t\_comp}), \text{t\_swap})\big) \]
\end{definition}

\begin{definition}[Denotational Mapping]
\lean{ISAR.denot}
\uses{ISAR.ExtTensor, ISAR.t_sₛ}
The structural denotational homomorphism mapping symbolic terms to the tensor space:
$$\text{denot} : \text{ITerm} \to \text{TensorSpace}$$
\end{definition}

\begin{theorem}[Denotational Soundness]
\lean{ISAR.denot_sound}
\uses{ISAR.denot}
One-step reduction in the symbolic calculus maps to extensional equivalence of denotations:
$$\forall t, u, \quad t \to_I u \implies \text{denot } t \approx_{\text{ext}} \text{denot } u$$
Thus, the denotational mapping factors uniquely through the Invariant Layer quotient:
$$\text{InvariantLayer.toExtTensor} : \text{InvariantLayer} \to \text{ExtTensor}$$
\end{theorem}

\subsection{Sparse Matrix Operator Composition}
In the concrete implementation of the axiomatic kernel (\url{scratch/isar_ski_kernel.py}), the emergent operator combinators are represented as products of the basic sparse matrices $I, R, A, S$:
\begin{align*}
\text{NORM} &\triangleq S \cdot I \\
\text{APP} &\triangleq A \cdot R \\
\text{DUP} &\triangleq S \cdot S \\
\text{COMP} &\triangleq R \cdot A \cdot S \\
\text{SWAP} &\triangleq R \cdot S \\
\text{CONST} &\triangleq I \cdot S
\end{align*}

\subsection{Plex-Agent Tensor Primitives and Special Forms}
Alternatively, the runtime environment and compiler stack defined in \url{tensor_primitives.py} specify four basic tensor primitives and four special forms derived from them:
\begin{align*}
\text{ATOM (0)} &\triangleq I \\
\text{PAIR (1)} &\triangleq S \\
\text{JOIN (2)} &\triangleq R \cdot A \\
\text{NORM (3)} &\triangleq I \cdot R \cdot A \cdot S
\end{align*}
The special forms are defined as compositions of these primitives:
\begin{align*}
\text{LAMBDA} &\triangleq \text{JOIN} \cdot \text{PAIR} \cdot \text{NORM} \\
\text{QUOTE} &\triangleq \text{ATOM} \cdot \text{NORM}
\end{align*}
along with conditional branching (\texttt{IF}) and variable bindings (\texttt{DEF}).

\section{Continuous Universal Approximation}

We formalize the continuous parameter space of the ISAR update.

\begin{definition}[Non-Polynomial Activation]
\lean{ISAR.Activation.nonPolynomial}
An activation function $\sigma \in C(\mathbb{R}, \mathbb{R})$ is non-polynomial:
$$\forall P \in \text{Polynomial } \mathbb{R}, \quad \sigma \neq P$$
\end{definition}

\begin{definition}[Raw Address Space]
\lean{ISAR.RawAddress}
The raw parameter configuration space $\text{RawAddress } d\ k$ representing a neural network configuration:
\begin{align*}
\text{RawAddress } d\ k &\triangleq \Sigma (N, T : \mathbb{N}), \quad (\text{Fin } T \to \text{Fin } 4 \to \mathbb{R}) \times C(\mathbb{R}^d, \text{GridState } N) \\
&\quad \times C(\text{GridState } N, \mathbb{R}^k)
\end{align*}
\end{definition}

\begin{definition}[Realization Map]
\lean{ISAR.realizeRaw}
\uses{ISAR.RawAddress}
The continuous function realized by a raw address configuration:
$$\text{realizeRaw } d\ k\ \sigma : \text{RawAddress } d\ k \to C(\mathbb{R}^d, \mathbb{R}^k)$$
\end{definition}

\begin{definition}[Address Equivalence]
\lean{ISAR.AddressEq}
\uses{ISAR.realizeRaw}
Two raw parameter configurations are equivalent if they realize the same continuous function.
\end{definition}

\begin{definition}[Kernel Address Space]
\lean{ISAR.KernelAddress}
\uses{ISAR.AddressEq}
The quotient parameter space of raw addresses modulo observational equivalence:
$$\text{KernelAddress } d\ k\ \sigma := \text{Quotient } (\text{addressSetoid } d\ k\ \sigma)$$
\end{definition}

We state the Universal Approximation Theorem (UAT).

\begin{axiom}[ISAR Universal Approximation]
\lean{ISAR.ISAR_UAT}
\uses{ISAR.Activation.nonPolynomial, ISAR.RawAddress, ISAR.realizeRaw}
Let $K \subset \mathbb{R}^d$ be compact, $f \in C(\mathbb{R}^d, \mathbb{R}^k)$ a target continuous function, $\sigma$ a non-polynomial activation, and $\varepsilon > 0$. Then there exists a raw address $\theta \in \text{RawAddress } d\ k$ such that:
$$\forall x \in K, \quad \| \text{realizeRaw } d\ k\ \sigma\ \theta\ x - f(x) \| < \varepsilon$$
\end{axiom}

\section{Quotient Completion and Unique Representation}

We complete the parameter space and show unique representation.

\begin{axiom}[Quotient Completion Space]
\lean{ISAR.KernelAddressLimit}
The metric completion of \texttt{KernelAddress} is \texttt{KernelAddressLimit}.
\end{axiom}

\begin{axiom}[Continuous Realization Limit]
\lean{ISAR.continuousRealizationLimit}
\uses{ISAR.KernelAddressLimit}
The extended continuous realization map on the completed space.
\end{axiom}

\begin{axiom}[Completion Embedding]
\lean{ISAR.kernelAddressEmbedding}
The canonical inclusion map $i : \text{KernelAddress} \to \text{KernelAddressLimit}$.
\end{axiom}

\begin{theorem}[Embedding Injectivity]
\lean{ISAR.kernelAddressEmbedding_injective}
\uses{ISAR.kernelAddressEmbedding}
The completion embedding is injective:
$$\forall q_1, q_2, \quad i(q_1) = i(q_2) \implies q_1 = q_2$$
\end{theorem}

\begin{theorem}[Embedding Density]
\lean{ISAR.kernelAddressEmbedding_dense}
\uses{ISAR.kernelAddressEmbedding}
The completion embedding has dense range (with respect to uniform convergence on compact domains).
\end{theorem}

\begin{axiom}[Topological Extension Bijection]
\lean{ISAR.topological_extension_bijection}
An injective map with a dense range into a complete metric space extends uniquely to a bijection on the completion of its domain.
\end{axiom}

\begin{theorem}[ISAR Representation Theorem]
\lean{ISAR.ISAR_representation}
\uses{ISAR.topological_extension_bijection, ISAR.kernelAddressEmbedding_injective, ISAR.kernelAddressEmbedding_dense}
Every continuous function $f \in C(\mathbb{R}^d, \mathbb{R}^k)$ has a unique address $\theta_{\text{limit}} \in \text{KernelAddressLimit}$ such that:
$$\text{continuousRealizationLimit } d\ k\ \sigma\ \theta_{\text{limit}} = f$$
\end{theorem}

\begin{theorem}[Physical System Address]
\lean{ISAR.physical_system_has_address}
\uses{ISAR.ISAR_representation}
Every physical system represented as an admissible continuous kernel has a unique address in the continuous completion limit space.
\end{theorem}

\chapter{Applications and Advanced Properties}

In this final chapter, we formalize several advanced computational and structural properties of the ISAR framework: Futamura projections for partial evaluation, observational isomorphisms between dialects, and the classification of referentially open systems.

\section{Partial Evaluation and Futamura Projections}

We formalize partial evaluation and compile-time optimization via Futamura projections.

\begin{definition}[Specializer]
\lean{ISAR.specialize}
The meta-level specializer (partial evaluator) substituting static variables:
$$\text{specialize} : \text{ITerm} \to (\mathbb{N} \to \text{Option } \text{ITerm}) \to \text{ITerm}$$
\end{definition}

\begin{definition}[PE Setup]
\lean{ISAR.PESetup}
\uses{ISAR.specialize}
Object-level mix setup: an evaluator, a meta specializer, its reflection as an object-language term \texttt{specTerm}, the mix equation, and \texttt{selfApp} (evaluating \texttt{specTerm} implements \texttt{spec}). Mix+selfApp alone admit a trivial residualizer; \texttt{Nontrivial} is a separate obligation. Jones--Gomard--Sestoft 1993 polyvariant offline mix remains open.
\end{definition}

We prove the three Futamura projections in this mix formulation.

\begin{theorem}[First Futamura Projection]
\lean{ISAR.futamura_first}
\uses{ISAR.specialize}
Mix equation at the substitution layer: specializing then substituting dynamic data equals substituting the full environment, under a coherence hypothesis on static/dynamic environments.
\end{theorem}

\begin{theorem}[Second Futamura Projection]
\lean{ISAR.futamura_second}
\uses{ISAR.PESetup}
Mix instantiation: evaluating the specialized specializer on source data yields the specialized interpreter. Holds for any \texttt{PESetup}; does not by itself give a compiler-generator of Jones--Gomard--Sestoft quality.
\end{theorem}

\begin{theorem}[Third Futamura Projection]
\lean{ISAR.futamura_third}
\uses{ISAR.PESetup}
Self-application mix instantiation: evaluating \texttt{specTerm} specialized to itself yields a compiler generator \emph{at the mix-equation level}. Online \texttt{JGS\_PE} covers the \texttt{IStep} signature; 1993 polyvariant cogen remains open.
\end{theorem}

\subsection{System Generation vs. Compilation}
It is important to distinguish the ontological role of the ISAR kernel from the epistemological role of the Futamura projections.
- \textbf{ISAR Kernel ($U(K) = I \cdot R \cdot A \cdot S$)}: Acts as a **System Generator**. It creates the computational substrate itself, establishing the tensor manifold in which computation exists and from which operators (norm, app, composition) emerge constructively.
- \textbf{Futamura Kernel ($I \times S \times A$)}: Acts as a **Compiler Generator**. It operates within an established space to transform program representations (interpreter $\to$ compiler).

\section{Dialect Isomorphisms and Unification}

We formalize structural unifications and observational equivalence between different dialects.

\begin{definition}[Observational Isomorphism]
\lean{ISAR.ObservationalIsomorphism}
An observational isomorphism between two dialects $D_1$ and $D_2$ is a pair of maps between their objects preserving the evaluation and coding semantics.
\end{definition}

\begin{definition}[Admissible Dialect]
\lean{ISAR.AdmissibleDialect}
A dialect packaged with its admissibility proofs, allowing it to induce a semantic kernel.
\end{definition}

\begin{theorem}[Isomorphism Unification]
\lean{ISAR.isomorphism_unification}
\uses{ISAR.AdmissibleDialect, ISAR.ObservationalIsomorphism}
An observational isomorphism between two admissible dialects induces an isomorphism between their corresponding semantic kernels in the category of kernels.
\end{theorem}

\section{Referential Openness and Anchor Dependency}

We analyze systems that are operationally closed versus referentially open (anchor-dependent).

\begin{definition}[Transition System]
\lean{ISAR.TransitionSystem}
A standard operationally closed autonomous transition system.
\end{definition}

\begin{theorem}[Forward Invariance]
\lean{ISAR.closure_preserved_under_reachability}
\uses{ISAR.TransitionSystem}
If a state-space subset $C$ is closed under the step relation (forward invariant), any state reachable from $C$ remains in $C$.
\end{theorem}

This applies to closed computational models (combinator engines, lambda calculus, stack VMs, set theory interpreters) where the state alone is sufficient to reconstruct the semantics at any step.

\begin{definition}[Anchor-Dependent System]
\lean{ISAR.AnchorDependentSystem}
A referentially open transition system whose transitions depend on an external environment context (anchor).
\end{definition}

\begin{theorem}[Anchor Dependency]
\lean{ISAR.referentially_open_requires_anchor}
\uses{ISAR.AnchorDependentSystem}
If a state can lead to different observations under different anchor sequences, the semantics cannot be resolved from the state alone without environmental anchors.
\end{theorem}

This establishes the boundary of decodability (the ``Reverse Rosetta'' problem). For open systems like natural language or undeciphered scripts (e.g., Linear A), analyzing syntax is insufficient to decode meaning because the external semantic anchors are lost. Thus, referentially open systems cannot be reconstructed from syntactic relationships alone.



\end{document}
